\section{APOIO AO DISCENTE}\label{APOIO}

Para o apoio aos discentes do \NOMECURSO, o Campus Maracanaú dispõe, além da coordenação do curso, de outros setores que promovem ações de orientação e acompanhamento pedagógico e psicológico, assim como serviços de assistência social, enfermagem e nutrição e núcleos de inclusão e acessibilidade.

O envolvimento, a participação e a colaboração dos setores como Coordenadoria de Assuntos Estudantis (serviço social, psicologia, enfermagem e nutrição), Coordenadoria de Controle Acadêmico, Coordenadoria Técnico-Pedagógica, Estágio, Biblioteca, entre outros, que também lidam com o corpo discente do campus, colaboram com a redução contínua da evasão e da retenção acadêmica, permitindo assim, que o estudante permaneça na instituição e conclua seu curso com êxito. As ações realizadas por cada setor ou serviço estão listadas a seguir:

%\clearpage
\subsection{Coordenação de Curso}
No que tange ao apoio discente, a Coordenação do \NOMECURSO visa ser facilitadora nas ações acadêmicas relacionadas ao curso e na resolução de possíveis demandas pelos alunos. Para tanto, lança mão de ações sistematizadas que vão desde o atendimento aos discentes, à promoção de estratégias para melhoria de processo de aprendizagem. 

Ademais, a coordenação do curso administra a potencialidade do corpo docente do seu curso, favorecendo a integração e a melhoria contínua. A atuação da coordenação do \NOMECURSO é detalhada em um plano de ação semestral ou anual elaborado com base nas diretrizes da Nota Técnica 4/2018/PROEN/REITORIA. O referido plano dispõe de indicadores de desempenho disponíveis e públicos e, de forma geral, esta é feita através de:

\begin{itemize}
	\item Resolver assuntos ligados ao aproveitamento de disciplinas e à entrada de alunos no curso por meio do Edital de Transferidos / Graduados;
	\item Tratar de assuntos pedagógicos do curso junto a Coordenação Pedagógica;
	\item Tratar de questões ligadas à matricula e situação regular dos alunos;
	\item Acompanhar a vida acadêmica do corpo discente;
	\item Coordenar as atividades relacionadas ao reconhecimento do curso;
	\item Solicitar dos professores os Programas de Unidades Didáticas (PUD) de todas as disciplinas do curso, bem como mantê-los atualizados;
	\item Atuar juntamente com o NDE e o Colegiado nas discussões sobre alterações na matriz curricular, quando se fizer necessário;
	\item Atuar juntamente com o NDE e o Colegiado na atualização do projeto pedagógico do curso, quando necessário;
	%\item Coordenar as atividades desenvolvidas pelos monitores e pelos bolsistas de laboratórios;
	\item Auxiliar ao setor administrativo em assuntos estratégicos, tais como: levantamento de demandas (infra-estrutura, equipamentos, etc.), elaboração de planos de trabalho, elaboração dos horários semestrais, levantamento de demandas de perfis de vagas para novos professores, planejamentos anuais, dentre outros.
	\item Liderar as ações de divulgação do curso na sociedade;
	\item Representar o curso em eventos e reuniões internas e externas, quando for o caso.
	\item Acompanhar o registro de aulas no acadêmico;
	\item Registrar aulas extras no acadêmico em comum acordo entre os professores e os discentes;
	\item Atender às solicitações de reabertura de diários no acadêmico em comum acordo entre professores e discentes;
	\item Dentre outras atividades detalhadas no plano de ação anteriormente mencionado.
\end{itemize}

\subsection{Coordenadoria de Controle Acadêmico}

A Coordenação de Controle Acadêmico (CCA), como órgão de execução, responde pelas questões operacionais junto ao Sistema Q-Acadêmico. Desse modo, define junto a Diretoria de Ensino/DIREN, a qual é subordinada, a execução dos processos de pré-matrícula, matrícula, criação de turmas e horários.

Pelo princípio da legalidade, executa procedimento em acordo com o Regulamento da Organização Didática/ROD, o que possibilita auxiliar coordenadores e estudantes quanto às diretrizes estabelecidos no regulamento, além de gerenciar procedimento de ingresso através do Sistema de Seleção Unificada – Sisu, através do acesso ao SISUGestão, o setor também controla e organiza arquivos de discentes. 

No atendimento ao público discente, emite documentação de situação acadêmica, como históricos, declarações e ementas das disciplinas aprovadas.

\subsection{Departamento de Extensão, Pesquisa, Pós-Graduação e Inovação}

%A Coordenadoria de Pesquisa, Pós-­Graduação e Inovação é um setor diretamente subordinado ao departamento de Extensão, Pesquisa, Pós-Graduação e Inovação – DEPPI, responsável por atividades de atendimento ao discente no que tange à concessão de auxílio acadêmico, auxílio didático-­pedagógico, bem como suporte aos discentes que participam de projetos de pesquisa como bolsistas com fomento ou bolsistas voluntários. Os auxílios são destinados aos alunos que participam de eventos científicos e/ou tecnológicos de âmbito nacional ou internacional.
%Em relação ao fomento da pesquisa, a coordenadoria divulga informações sobre editais internos e externos, além de dar orientação sobre os procedimentos de cadastro de bolsistas e acompanhar o andamento das atividades de pesquisa. Ademais, os discentes participantes do programa de Iniciação Científica recebem apoio para divulgar os resultados de suas pesquisas através do Seminário de Iniciação Científica e Tecnológica (SEMIC).

O Departamento de Extensão, Pesquisa, Pós-Graduação e Inovação (DEPPI) do  Campus Maracanaú é responsável por atividades de atendimento ao discente no que tange a concessão de auxílio acadêmico, auxílio didático-pedagógico, bem como suporte aos discentes que participam de projetos de extensão, pesquisa e inovação como bolsistas com fomento ou sem fomento por meio do auxílio acadêmico. O departamento apoia os discentes que aprovam trabalhos em eventos científicos e/ou tecnológicos de âmbito nacional ou internacional. 

Em relação ao fomento da extensão, pesquisa e inovação, o DEPPI divulga informações sobre editais internos e externos, além de dar orientação sobre os procedimentos de cadastro de bolsistas e acompanhar o andamento das atividades de pesquisa e dos projetos de extensão. 

Esse departamento, ainda, realiza captação de vagas de estágio por meio de seu programa de Relações Empresariais, no qual mantém estreita relações com as empresas locais e trabalha em conjunto com o Setor de Estágios e a Diretoria de Ensino (DIREN). 

Ademais, o DEPPI organiza eventos de extensão e pesquisa como a Semana de Integração Científica (SIC) e o Seminário de Iniciação Científica e Tecnológica (SEMIC). Nesses eventos, os discentes participantes dos programas de Iniciação Científica e de Extensão recebem apoio para divulgar os resultados de seus trabalhos para a comunidade.

\subsection{Coordenadoria Técnico Pedagógica}

A Coordenadoria Técnico-Pedagógica (CTP) do IFCE é o setor responsável pelo planejamento, acompanhamento, avaliação de ações pedagógicas desenvolvidas no campus com vistas à formulação e reformulação contínua de intervenções pedagógicas que favoreçam o alcance de resultados satisfatórios quanto ao processo ensino-aprendizagem.

A atuação da CTP é embasada nos fundamentos e pressupostos teóricos educacionais, nos princípios legais da Educação Brasileira. A atuação desse setor encontra-se em consonância com a Missão Institucional do IFCE. A seguir, apresentam-se as atribuições do referido setor:

\begin{itemize}
	\item Realizar atendimento individual e/ou em grupo aos estudantes, professores, pais e ou responsáveis conforme necessidades observadas pelo setor e ou quando solicitado;
	\item Analisar continuamente as causas da evasão e repetência para formular sistematicamente estratégias que visem à superação ou minimização do problema;
	\item Acompanhar o desenvolvimento dos estudantes com baixo rendimento escolar (frequência e desempenho) propondo alternativas que favoreçam a superação e a minimização dessa problemática;
	\item Mediar à relação professor e aluno e aluno - aluno voltada para o sucesso do desempenho discente solicitando apoio e parceria da Assistência Estudantil e Setor de Psicologia, quando necessário;
	\item Realizar a cada período letivo, a Avaliação de Desempenho Docente, com vistas a promover a melhoria da prática docente por meio de análise dos resultados da avaliação e dos feedbacks que serão dados por meio de conversas individuais e capacitações;
	\item Promover ações formativas (encontros pedagógicos, encontros de estudo, capacitações, orientações individuais, conselhos de classe, colegiados, entre outros) que provoquem no professor avaliação de sua prática docente para que ao longo do processo melhorem sua atuação pedagógica.
\end{itemize}

Convém destacar que as atribuições da CTP se articulam com as ações desenvolvidas por outros setores da instituição, como coordenações de cursos, coordenadoria de assistência estudantil (serviço social, psicologia, enfermagem e nutrição), coordenadoria de controle acadêmico, estágio, biblioteca, pesquisa, extensão, entre outros, que também lidam com o corpo discente do campus.

\subsection{Coordenadoria de Assuntos Estudantis}

A Assistência estudantil vem se consolidando no Instituto Federal de Educação, Ciência e Tecnologia do Ceará – IFCE como um conjunto de ações, configurando-se através de auxílios financeiros e serviços, visando ampliar as condições de permanência e apoio à formação acadêmica do corpo discente. Uma dessas ações diz respeito à disponibilização de serviços, caracterizados por ações continuadas, visando ao atendimento biopsicosocial do discente. Outra ação diz respeito aos auxílios sob a forma de pecúnia, sendo estes destinados, na sua maioria, ao discente, prioritariamente em condições de vulnerabilidade social, e operacionalizados por meio do regulamento dos auxílios. Tal regulamento é normatizado pelo programa de Auxílios, previsto na Política de Assistência Estudantil do IFCE ( aprovada pela resolução nº 024, de 22 de julho de 2015) e, institui ações de efetivação do Decreto nº 7.234, de 19 de junho de 2010, que dispõe sobre o Programa Nacional de Assistência Estudantil (PNAES).

As ações previstas na PNAES dizem respeito às seguintes áreas: moradia estudantil, alimentação, transporte, atenção à saúde, inclusão digital, cultura, esporte, creche, apoio pedagógico, acesso e participação e aprendizagem de estudantes com deficiência, transtornos globais do desenvolvimento e altas habilidades e superdotação (Decreto 7.234/2010, Art. 3º). Ressaltamos, ainda, que o referido decreto prevê que estas ações serão executadas por Instituições Federais de Ensino Superior, contemplando os IFs. 
Portanto, a assistência Estudantil no IFCE, vislumbrada mediante serviços ofertados (merenda escolar, atendimento psicológico, atendimento pedagógico, entre outras ações) e auxílios financeiros foram instituídos na perspectiva de “viabilizar a igualdade de oportunidades, contribuir para a melhoria do desempenho acadêmico e agir, preventivamente, nas situações de retenção e evasão decorrentes da insuficiência de condições financeiras.

O IFCE Campus Maracanaú dispões dos seguintes serviços, diretamente subordinados à Coordenadoria de Assuntos Estudantis, a saber: Serviço de Enfermagem, Serviço de Nutrição, Serviço de Psicologia e Serviço Social. As ações realizadas por cada serviço estão listas a seguir:

\subsubsection{Serviço de Enfermagem} 

No âmbito do IFCE, a Enfermagem destina-se a promoção da saúde com foco na educação em saúde, bem como a oferecer cuidados de primeiros socorros em situações de urgência e emergência, conforme ações elencadas a seguir: 

\begin{itemize}
	\item contribuir para o desenvolvimento integral do(a) discente; 
	\item colaborar no mapeamento da realidade socioeconômica, acadêmica e de saúde dos discentes; 
	\item apoiar as estratégias de inclusão das pessoas com deficiência;
	\item atuar na prevenção, promoção, tratamento e vigilância à saúde de forma individual e coletiva, colaborando com o processo de ensino-aprendizagem; 
	\item realizar ações de prevenção e controle sistemático de situações de saúde e agravos em geral; 
	\item desenvolver atividades de educação em saúde para a adoção de hábitos saudáveis, visando à melhoria da qualidade de vida e à promoção da saúde da comunidade acadêmica; 
	\item participar de estratégias de combate à evasão escolar; 
	\item participar do planejamento, execução e avaliação da programação das ações anuais de saúde; 
	\item participar do processo de seleção de auxílios referente aos aspectos relativos às situações de saúde;
	\item acompanhamento de discentes aos serviços de saúde, nas situações previstas nas diretrizes para atuação do enfermeiro no IFCE;
	\item programa Saúde e Prevenção nas Escolas(SPE);
	\item camisinha Card;
	\item vacinação, sendo previstas a realização de 06 campanhas anuais;
	\item campanhas semestrais de doação de sangue;
	\item atendimento ambulatorial: aconselhamento em DST/HIV e AIDS; realização de curativos; aferição de pressão arterial; glicemia; auscultas cardíacas e sinais vitais; primeiros socorros; atendimento aos servidores e alunos com hipertensão e
	diabetes.
\end{itemize}

\subsubsection{Serviço de Nutrição} 

O Serviço de Alimentação e Nutrição é responsável pela administração da Unidade de Alimentação no campus, a qual visa à oferta de uma alimentação adequada, compreendendo o uso de alimentos variados, seguros, que respeitem a cultura, as tradições e os hábitos alimentares saudáveis, contribuindo, assim, para melhoria do rendimento escolar, permanência do estudante no espaço educacional e promoção de hábitos alimentares saudáveis.

O Serviço de Nutrição ainda atua nos programas de educação e assistência nutricional, desenvolvendo ações com a equipe multiprofissional tendo em vista a promoção da saúde e segurança alimentar e nutricional, prestando, também, assessoria às atividades de ensino, pesquisa e extensão.


\subsubsection{Serviço de Psicologia} 

A psicologia escolar/educacional assume um papel de contribuir para a construção de uma educação de qualidade, baseada nos princípios do compromisso social, do respeito à diversidade e dos direitos humanos. Entende que a ação educativa é permeada por determinantes biopsicossociais que interferem, direta e indiretamente, no desenvolvimento do processo de aprendizagem de cada individuo, desse modo a ação educativa não se limitará a queixa, mas a busca constante de fomentar um ambiente escolar que promova saúde mental.

Neste sentido, o serviço de Psicologia do IFCE - Campus Maracanaú busca:

\begin{itemize}
	\item Apoiar servidores no trabalho com a heterogeneidade de discentes;
	\item Avaliar, acompanhar e orientar dentro do contexto institucional casos que requeiram encaminhamentos clínicos, estabelecendo um espaço de acolhimento, escuta e reflexão. No caso de demandas psicoterápicas, será realizado encaminhamento para outras instituições que ofereçam o tratamento adequado;
	\item Fazer parte da equipe multiprofissional que envolve o processo de ensino e aprendizagem levando em conta o desenvolvimento global do discente;
	\item Propiciar condições para que o discente expresse sua autonomia e consciência crítica, por meio da participação ativa na vida acadêmica, contribuindo para uma formação cidadã; 
	\item Realizar acompanhamento dos discentes em situação de vulnerabilidade socioeconômica e dificuldade de aprendizagem para a realização das intervenções necessárias;
	\item Identificar e analisar as causas e as motivações das reprovações, retenções e evasões dos discentes, a fim de subsidiar o direcionamento das intervenções, apreendendo quais os aspectos sociais, físicos, cognitivos e afetivos geram resistência no seu processo de aprendizagem elaborando condições para permanência da qualidade da aprendizagem;
	\item Propiciar aos discentes espaços de reflexão e diálogo sobre as temáticas demandadas pelos diversos atores que compõem a comunidade acadêmica; 
	\item Fomentar momentos de expressões artísticas, espirituais, culturais e esportivas do discente e comunidade acadêmica, propiciando as inter-relações e a circulação da palavra nas suas mais diferentes manifestações;
	\item Estimular a criatividade e iniciativa dos discentes para criação de grupos autogeridos que trabalhem temáticas por eles definidas;
	\item Favorecer a prevenção e promoção da saúde dos discentes e comunidade acadêmica, visando o alcance da discussão dos diversos aspectos que compõem o conceito ampliado de saúde, a partir de trabalhos preventivos que visem um processo de transformação pessoal e social;
	\item Promover ações articuladas com a rede socioassistencial, educacional e de saúde do município, inserindo o campus Maracanaú como um dos pontos estratégicos de mobilização social do município.
\end{itemize}


\subsubsection{Serviço Social}

O Serviço Social no Campus de Maracanaú insere-se na promoção do Programa Nacional de Assistência Estudantil (PNAES – Decreto MEC Nº 7234), mediante elaboração e implementação de serviços, programas, projetos e auxílios (sob a forma de pecúnia), visando à ampliação das condições de acesso e de permanência, com enfoque numa formação crítica e autônoma.

A atuação do Serviço Social no Campus situa-se no âmbito da Assistência Estudantil, com destaque nas seguintes ações:

\begin{itemize}
	\item \textbf{De caráter individual:} atendimento social, escuta qualificada, estudo social, análise socioeconômica, socialização de informações, orientações sociais, encaminhamento para outros serviços, seleção de estudantes para concessão de auxílios.
	\item \textbf{De caráter coletivo:} atendimento coletivo, formação de grupos, reuniões, encontros, seminários, oficinas para alunos e técnicos, campanhas, realização de atividades de acolhimento e integração dos discentes à comunidade acadêmica, confecção de materiais educativos, mobilização e organização social e política, apoio à constituição das entidades estudantis, capacitação dos alunos e técnicos, participação nos espaços de controle social.
\end{itemize}

Destacamos que é de responsabilidade do Serviço Social, a concessão dos auxílios financeiros, a saber: 

\begin{itemize}
	\item AUXÍLIO MORADIA - subsidia despesas com habitação para locação, sublocação de imóveis para discentes com referência familiar e residência domiciliar fora da Sede do município onde está instalado o campus;
	\item AUXÍLIO ALIMENTAÇÃO - subsidia despesas de alimentação nos dias letivos;
	\item AUXÍLIO TRANSPORTE – subsidia despesas no trajeto residência/campus/residência;
	\item AUXÍLIO ÓCULOS – complementa despesas de aquisição de óculos ou lentes corretivas de deficiências oculares;
	\item AUXÍLIO VISITAS/VIAGENS TÉCNICAS – subsidia despesas com alimentação e/ou hospedagem, em visitas e viagens técnicas;
	\item AUXÍLIO ACADÊMICO – complementa despesas com alimentação, hospedagem, passagem e inscrição dos discentes para a participação em eventos acadêmicos;
	\item AUXÍLIO DIDÁTICO-PEDAGÓGICO – subsidia a aquisição de material de uso individual e intransferível, indispensável à aprendizagem de determinada disciplina;
	\item AUXÍLIO DISCENTES MÃES/PAIS – subsidia despesas de filho(s) de até 06 (seis) anos de idade ou com deficiência, sob sua guarda;
	\item AUXÍLIO FORMAÇÃO – subsidia despesas relativas à ampliação da formação dos discentes em laboratórios/oficinas e em projetos caracterizados por ensino, pesquisa e extensão, vinculados ao seu curso.
\end{itemize}

Os auxílios têm por objetivos e finalidades ampliar as condições de permanência e apoio à formação acadêmica dos discentes, visando a reduzir os efeitos das desigualdades sociais; contribuir para reduzir a evasão; propiciar a melhoria do desenvolvimento acadêmico e biopsicossocial do discente.

\subsection{Biblioteca}

A Biblioteca Rachel de Queiroz, fundada em 12 de março de 2010, ocupa uma área de 390,11 m2, possui 62 assentos para estudo individual ou em grupo e oferece suporte para o ensino, pesquisa e extensão. Conta com um acervo de 1.850 títulos e 11.900 exemplares entre livros, cd-rom, disquete e folhetos, nas áreas técnicas, tecnólogos, licenciatura e engenharia e tecnologia, com ênfase em livros técnicos e didáticos. Conta ainda com 88 títulos e 407 exemplares de revistas doadas, além de acesso ao portal de Periódicos da CAPES\footnote{www.periodicos.capes.gov.br}, disponível para todos os computadores do IFCE Maracanaú.

Atualmente a Biblioteca Rachel de Queiroz disponibiliza a Biblioteca Virtual Universitária \footnote{https://ifce.bv3.digitalpages.com.br/login} contemplando todo o acervo bibliográfico da lista de títulos da Pearson e das editoras parceiras: Manole, Contexto, Ibpex, Papirus, Casa do Psicólogo, Ática, Scipione, MartinsFontes, Cia das Letras, Rideel, Educs e Jaypee. Todo acervo destas editoras pode ser acessado, de forma digital, por professores e alunos do IFCE. Dentre as muitas áreas beneficiadas destacam-se Ciências da Computação, Engenharias, Física, Matemática e Estatística. Além disso, a biblioteca presta serviços como o empréstimo domiciliar de todos os materiais que compõem o acervo; a consulta à base de dados tanto nos terminais de autoatendimento local quanto via internet; o acesso ao Portal de Periódicos Eletrônicos da Capes; a elaboração de catalogação na fonte; a orientação técnica para elaboração e apresentação de trabalhos acadêmicos, com base nas normas técnicas de documentação da ABNT; e levantamentos bibliográficos e referenciais para pesquisas.

A biblioteca conta com profissionais que registram e catalogam, classificam e indexam as novas aquisições e fazem a manutenção das informações bibliográficas no Sistema Sophia além de realizar preparação física (carimbos de identificação e registro, colocação de etiquetas) do material bibliográfico para empréstimo domiciliar.

A Biblioteca funciona de segunda a sexta-feira, nos seguintes horários:

\begin{itemize}
	\item Empréstimo, estudo e leitura - 8h às 20h
	\item Referência (Consulta Local) e Periódicos - 8h às 20h
	\item Multimídia (computadores conectados à Internet) - 8h às 20h.
\end{itemize}

%
%A biblioteca possui um acervo com mais de 3.700 volumes, entre livros e CD’s, e mais 250 revistas periódicos nas áreas
%de ciências humanas, ciências puras, artes, literatura e tecnologia, com ênfase em livros técnicos e didáticos. São 904
%títulos nas mais diversas áreas do conhecimento.
%
%A biblioteca conta com profissionais que registram e catalogam, classificam e indexam as novas aquisições e fazem a
%manutenção das informações bibliográficas Sistema GNUTECA. Realizam, também, a preparação física (carimbos de
%identificação e registro, colocação de etiquetas, bolso e fichas de empréstimo) do material bibliográfico para
%empréstimo domiciliar.
%
%Principais serviços:
%
%\begin{itemize}
%\item Acesso à Base de Dados Gnuteca nos terminais locais e via Internet.
%\item Empréstimo domiciliar e renovação das obras e outros materiais.
%\item Consulta local ao acervo.
%\item Elaboração de catalogação na fonte.
%\item Orientação técnica para elaboração e apresentação de trabalhos acadêmicos, com base nas Normas Técnicas de
%Documentação da ABNT.
%\item Acesso ao Portal de Periódicos da Capes.
%\item Acesso à Internet.
%\item Levantamento bibliográfico.
%\end{itemize}
%O \textit{campus} conta ainda com um serviço adicional da Biblioteca Virtual, acessível via internet para todos os alunos
%matriculados. Este serviço disponibiliza um grande número de títulos para acesso por computador, dentro ou fora da
%insituição.

\subsubsection{Empréstimo}
O usuário poderá retirar, por empréstimo domiciliar, qualquer publicação constante do acervo bibliográfico, exceto as
obras de referência (enciclopédias, dicionários, atlas, periódicos, jornais, etc) ou outras publicações que, a critério
da Biblioteca, não podem sair.

\subsubsection{Setor Multimídia}
O Setor de Multimídia possui 5 computadores conectados à Internet para que o usuário possa fazer pesquisas online.
Funciona de segunda a sexta-feira, das 8h às 20h.

\subsection{Núcleo de Acessibilidade às Pessoas com Necessidades Educacionais Específicas} 

O Núcleo de Acessibilidade às Pessoas com Necessidades Educacionais Específicas (NAPNE) tem como objetivo disseminar uma cultura da “educação para convivência”, aceitação da diversidade e, principalmente, buscar a quebra das barreiras arquitetônicas, tecnológicas, educacionais e atitudinais.

Para tanto o NAPNE atua no sentido de:

\begin{itemize}
	\item Buscar a quebra de barreiras arquitetônicas, comunicacionais, educacionais e atitudinais na Instituição de ensino, por meio de levantamentos e aplicação de questionários periodicamente;
	\item Promover condições necessárias para o ingresso, a permanência e o êxito educacional de discentes com necessidades educacionais específicas no IFCE, realizando o acompanhamento dos estudantes;
	\item Atuar junto às coordenações de cursos, à equipe pedagógica e aos colegiados dos cursos oferecendo suporte no processo de ensino-aprendizagem dos estudantes com necessidades educacionais específicas, colaborando com a adaptação dos referenciais teórico-metodológicos, colocando a equipe à disposição para prestar esclarecimentos e orientações;
	\item Articular junto ao Campus e à PROEXT a disponibilização de recursos específicos para aquisições de materiais de consumo e permanente que possibilitem a promoção das atividades de ensino, pesquisa e extensão com qualidade;
	\item Potencializar o processo ensino-aprendizagem por meio da utilização de novas tecnologias de informação e de comunicação (TICs) que facilitem esse processo, por meio da indicação dos recursos já existentes, assim como colaborando com projetos e pesquisas, e ainda promovendo campanha de conscientização e incentivo a ações inclusivas (Prêmio IFCE Inclusivo – premiação de honra ao mérito por ações, projetos e produtos desenvolvidos no IFCE Maracanaú);
	\item Promover e participar de estudos, eventos e debates sobre Educação Inclusiva com o intuito de informar e sensibilizar a comunidade acadêmica no âmbito do IFCE e de outras instituições, realizando palestras e rodas de debates (Projeto Encontros Inclusivos), além do curso de Libras (Módulos I, II e III, totalizando 120hs);
	\item Contribuir para a inserção da pessoa com necessidades educacionais específicas no IFCE e em espaços sociais, realizando a divulgação dos editais de seleção e dos cursos em instituições que atuem com pessoas com deficiência, além de fazer parceria com o Conselho Municipal dos Direitos da Pessoas com Deficiência de Maracanaú e Associações aproximando-os do campus.
	\item Assessorar a Diretoria de Ingressos do IFCE especificamente nos casos de ingresso de estudantes e servidores com necessidades específicas, formando uma comissão para o acompanhamento da análise dos documentos dos cotistas no processo de matrícula.
	\item Assessorar, quando necessário, no processo de alterações nas regulamentações que visem o ingresso e a permanência de pessoas com necessidades educacionais específicas no IFCE.
\end{itemize}

\subsection{Setor de Estágio} 

%O Setor de Estágio do campus de Maracanaú do Instituto Federal do Ceará é diretamente subordinado ao departamento de Extensão, Pesquisa, Pós-Graduação e Inovação – DEPPI e responsável pela administração do estágio discente, seja ele obrigatório ou não-obrigatório. Atua, em parceria, com a direção de ensino e coordenações de cursos, e conta com o apoio dos professores orientadores de estágio.
%Ainda, realiza o controle das documentações, acompanhamento dos relatórios e o cumprimento das regras de estágio conforme Lei n° 11.788, de 25 de Setembro de 2008.
%Também, faz a divulgação das ofertas de estágio pelas empresas para disseminar as oportunidades ao corpo discente.

O Setor de Estágio do IFCE-Campus de Maracanaú é diretamente subordinado à Diretoria de Ensino e é responsável pela administração do estágio discente, seja ele obrigatório ou não-obrigatório. Atua, em parceria, com o Departamento de Extensão, Pesquisa, Pós-Graduação e Inovação (DEPPI) e coordenações de cursos, e conta com o apoio dos docentes orientadores de estágio.

Ainda, realiza o controle das documentações, o acompanhamento dos relatórios e cumprimento das regras de estágio conforme Lei N° 11.788, de 25 de setembro de 2008, bem como a divulgação das ofertas de estágio pelas empresas para disseminar as oportunidades ao corpo discente.

\subsection{Setor de Educação Física e Esporte}

O Setor de Educação Física e Esporte (SEFE) oferece a toda a comunidade acadêmica do Campus Maracanaú além de uma avaliação física sistemática, diversas possibilidades para a prática de atividade física e esportes, entre elas: musculação, natação, hidroginástica, treinamento funcional, futebol de campo, futebol de salão, voleibol de quadra, voleibol de areia, futevôlei, basquetebol, handebol, tênis de mesa e jogos de tabuleiro.

O SEFE ainda possibilita ao público discente compor suas seleções esportivas e participar das competições a nível regional (jogos do IFCE sub-19 e aberto) e nacional (jogos dos IF sub-19). Além disso, possibilita também a socialização e integração entre discentes, docentes e comunidade por meio dos projetos de extensão desenvolvidos no setor.

% CORPO DOCENTE